\subsection{Kataoka's two odd rows with a repeated diagonal factor}

The first list contains a basis, and the second contains $(2,1),(1,1)$, another basis. Thus all layers are connected.

Put $d=\lambda-2$, $m=2^s$, $e=2^{s-1}$ and let $F$ be the Laplacian symbol. Direct multiplication gives
\[
 \overline F\doteq (X+1)(Y+1)(XY+1)^d
 \quad\text{in }\mathbb F_2[X^{\pm1},Y^{\pm1}],
\]
where $\doteq$ means equality up to a Laurent monomial. For each primitive pair $(a,b)$ modulo $m$, exactly one of $a,b,a+b$ is even. Write $c$ for this even member and $W=\zeta^c$, $\zeta$ primitive of order $m$. Let $2^t$ be the order of $W$, including $t=0$ for $W=1$. For $t\ge1$ the order of $(1+z)^c-1$ in characteristic two is $B=e/2^{t-1}$, taking $c$ to be its residue in $[0,m)$. Replacing an exponent by a congruent one does not change the formal order below $m$. At the axis $a$ even or $b$ even the contact order is $B+d+1$. At the diagonal $a+b$ even it is $2+dB$. If the contact order is strictly below $e$, it equals the normalized valuation $e v_2(F(\zeta^a,\zeta^b))$ by the uniformizer test. For $s\ge5$ this settles both axes when $t\ge2$, and the diagonal when $t\ge L+1$, where $L=2$ for $d=3$ and $L=3$ for $d=5$.

For the remaining classes let $h=(\lambda-1)/2$ and let $Z$ denote a primitive root of order $m$. In an axis class expand $W^hZ^hF(W,Z)$ as $\sum_j c_j(Z-1)^j$; symmetry treats the second axis. In a diagonal class use $W^h ZF(Z,W/Z)$ instead. The following table gives the smallest pair $(v_2(c_j),j)$ in lexicographic order. A dash means that the generic calculation applies.
\[
\begin{array}{c|cc|cc}
&\multicolumn{2}{c|}{\text{axis}}&\multicolumn{2}{c}{\text{diagonal}}\\
 t&\lambda=5&\lambda=7&\lambda=5&\lambda=7\\\hline
 0&(1,2)&(1,4)&(2,2)&(2,2)\\
 1&(1,2)&(1,4)&(3,0)&(4,0)\\
 2&-&-&(2,0)&(3/2,2)\\
 3&-&-&-&(2,0)
\end{array}
\]
To verify the entries, at $W=1,-1$ respectively the axis coefficient lists, in increasing $j$, are
\[
\begin{array}{c|ll}
5&(0,0,-10,-10,-2)&(8,16,6,-2,0)\\
7&(0,0,-28,-56,-42,-14,-2)&(-16,-48,-56,-32,-10,-2,0).
\end{array}
\]
For the diagonal the lists in increasing $t$ are
\[
\begin{array}{c|l}
5&(0,0,-4),\quad(8,8,0),\quad(-12,-12,0)\\
7&(0,0,-4),\quad(-16,-16,0),\quad
 (-16W,-16W-4,-2W-2),\quad(12W^3,12W^3,0).
\end{array}
\]
These expansions use only the identity $\Phi_{2^t}(W)=0$ and $\Phi_1(W)=W-1$ when $t=0$. For $W$ of order four we have $v_2(1+W)=1/2$, explaining the fractional entry. If the table entry is $(H,J)$, the normalized valuation is $eH+J$: the coefficient valuations lie in $\varphi(2^t)^{-1}\mathbb Z$. For $s\ge5$, the polynomial degree (at most six on the axes and two on the diagonal) is smaller than $e/\varphi(2^t)$ in every tabulated case. Thus the minimum term is unique.

To sum each class over its primitive pairs, for each fixed $W$ there are $e$ choices for the primitive coordinate $Z$. A normalized order $D$ therefore contributes $D$ for each such $W$, and there are $\varphi(2^t)$ choices of exact order $2^t$. In a single axis class the generic sum is $(s-2)e+(d+1)(e-2)$, while its two exceptional classes contribute $2(e+d-1)$. Each axis total is consequently $(s+d+1)e-4$. The generic diagonal sum is
\[
 d(s-L-1)e+2(e-2^L).
\]
Its exceptional classes contribute $9e+2$ for $d=3$ and $17e+6$ for $d=5$. Adding all three classes yields
\[
 W_s=(5s+10)e-14\quad (\lambda=5),\qquad
 W_s=(7s+11)e-18\quad (\lambda=7),\qquad s\ge5.
\]
The four initial valuations of the trees are $(8,34,134,358)$ and $(10,50,154,446)$ respectively. They can be verified by the finite cyclotomic basis calculation given above; the accompanying script recomputes all twelve initial layers. Since $T_s-T_{s-1}=W_s-2$, summation proves the asserted formulas.
